\ssr{Введение}

Цели данной лабораторной работы:
\begin{enumerate}
	\item изучить основные методы управления памятью в Linux;
	\item построить карту виртуальной памяти процесса;
	\item приобрести практические навыки работы с инструментами мониторинга использования памяти, настройки параметров виртуальной памяти, тестирования работоспособности оперативной памяти.
\end{enumerate}

\ssr{Выполнение}

\ssection{Задание 1. Получение карты виртуальной памяти процесса}

\begin{figure}[!htbp]
	\centering
	%\includegraphics[width=0.6\linewidth]{""}
	\caption{Полученная карта}
\end{figure}
\FloatBarrier

\ssection{Задание 2. Настройка файла подкачки}

\begin{figure}[!htbp]
	\centering
	%\includegraphics[width=0.6\linewidth]{""}
	\caption{Текущее использование памяти}
\end{figure}
\FloatBarrier

\begin{figure}[!htbp]
	\centering
	%\includegraphics[width=0.6\linewidth]{""}
	\caption{Создание и подключение файла подкачки}
\end{figure}
\FloatBarrier

\begin{figure}[!htbp]
	\centering
	%\includegraphics[width=0.6\linewidth]{""}
	\caption{Проверка}
\end{figure}
\FloatBarrier

\begin{figure}[!htbp]
	\centering
	%\includegraphics[width=0.6\linewidth]{""}
	\caption{Деактивация файла подкачки}
\end{figure}
\FloatBarrier

\ssection{Задание 3. Тестирование памяти}

\begin{figure}[!htbp]
	\centering
	%\includegraphics[width=0.6\linewidth]{""}
	\caption{Процесс и результаты тестирования}
\end{figure}
\FloatBarrier

\ssection{Дополнительное задание}

\begin{figure}[!htbp]
	\centering
	%\includegraphics[width=0.6\linewidth]{""}
	\caption{Значение параметра swappiness}
\end{figure}
\FloatBarrier

\ssr{Вывод}

В ходе лабораторной работы были:  
\begin{itemize}
	\item изучены основные методы управления памятью в Linux;
	\item построена карта виртуальной памяти процесса;
	\item приобретены практические навыки работы с инструментами мониторинга использования памяти, настройки параметров виртуальной памяти, тестирования работоспособности оперативной памяти.
\end{itemize}

\ssr{Контрольные вопросы}

\begin{enumerate}
	\item \textit{Какие задачи решает подсистема управления памятью?}\newline 
	
	\item \textit{Какой модуль занимается пересчетом виртуальных адресов в физические? Какие данные для работы модуля формирует ОС?}\newline 
	
	\item \textit{Что такое виртуальная память? Каковы преимущества использования ВАП по сравнению с работой с физическими адресами напрямую?}\newline 
	
	\item \textit{Какие адреса отображаются в карте ВАП?}\newline 
	
	\item \textit{Что такое файл подкачки? Как он используется?}\newline 
	
	\item \textit{Какая информация содержится в файле /etc/fstab?}\newline 
	
	\item \textit{Для чего используется утилита memtester?}\newline 
	
	\item \textit{Если при выполнении команды mkswap в качестве параметра передан файл размером 2Gb с разрешениями 777, что будет выведено в терминал?}\newline 
	
	\item \textit{В чем разница между файлом подкачки и разделом подкачки?}\newline 
	
	\item \textit{Как можно посмотреть количество свободной физической памяти? Приведите варианты команд.}\newline 
	
\end{enumerate}
